\documentclass{article}

\title{A Small Example for Writing-Complexity Analysis}
\author{J.C. Vaught}

\begin{document}
\maketitle

\begin{abstract}
Clear technical writing helps readers evaluate a result without unnecessary delay. This example contains ordinary prose and a few deliberately dense constructions for comparison.
\end{abstract}

\section{A Direct Explanation}

The sensor records pressure at a fixed rate. Each observation enters a compact table. The analysis compares the measured response with the expected trend.

\section{A Denser Explanation}

The calibration procedure, which was developed after several preliminary experiments had revealed temperature-dependent drift in the reference instrument, adjusts every observation before the estimator computes the final response.

This shows why a long interruption between a subject and its verb can slow a reader. The main target metric of this paper remains difficult to identify when the phrase does not name it directly.

\begin{equation}
  y = a x + b
\end{equation}

The displayed equation is excluded from the prose statistics. The concluding sentence remains available for analysis.

\begin{thebibliography}{1}
\bibitem{example} A reference title that should not enter the prose statistics.
\end{thebibliography}

\end{document}
